% Source: https://artofproblemsolving.com/community/c3h1410098_interesting_problem
% Stable identity: aops-topic-1410098-q01
In one revolution a point at distance $\rho$ from the centre travels once round a circle of circumference $2\pi\rho$.

\begin{asy}
usepackage("amsmath");
size(5cm);
defaultpen(fontsize(10pt));
real r1=3,r2=9;
fill(circle((0,0),r2),palegray); fill(circle((0,0),r1),white);
draw(circle((0,0),r2),linewidth(0.9)); draw(circle((0,0),r1),linewidth(0.9));
dot((0,0)); label("$O$",(0,0),SW);
pair P=r2*dir(30),Q=r1*dir(150);
draw((0,0)--P,gray+dashed); draw((0,0)--Q,gray+dashed);
dot(P,red); dot(Q,blue); label("$P$",P,dir(30),red); label("$Q$",Q,dir(150),blue);
draw(P--P+4.5*dir(120),red,Arrow(6)); draw(Q--Q+1.5*dir(240),blue,Arrow(6));
label("$3$",(0,0)--Q,dir(60)); label("$9$",(0.6*dir(30))--P,dir(120));
label("\mbox{}",(min(currentpicture,user=true)+max(currentpicture,user=true))/2,black); // keeps the text after the figure black when LaTeX typesets the labels inline
\end{asy}

At $20$ revolutions per minute:
\[
\text{speed of }P=20\cdot2\pi\cdot9=\boxed{360\pi\ \text{cm per minute}}\approx1131,\qquad \text{speed of }Q=20\cdot2\pi\cdot3=\boxed{120\pi\ \text{cm per minute}}\approx377 .
\]
Both points sweep out the same \emph{angle} per minute ($20$ full turns), so their angular speeds are equal. But in the same time $P$ covers a circle three times as long as $Q$ does, so its speed (distance per unit time) is three times as large. In general, speed $=$ (angular speed) $\times$ (distance from the centre): equal revolutions per minute means equal angle per minute, not equal distance per minute, so there is no contradiction.
